% answers_short.tex —— 2014 年·辽宁卷·文科 · 简短答案速查
% 由 tools/make_short_answers.py 自 answers.tex 的 \gksolhead/\gkanswer 提取，
% 并与 2014/2014_数学真题_辽宁文_答案要点.md 逐题比对后生成，勿手改（改源 MD 后重跑）。

\begin{gknotice}
\textbf{使用说明}：本节只给最终答案便于快速核对，完整推导见解析版。答案为编写者独立推导，\textbf{非官方标准答案}。
\end{gknotice}

\gkpart{选择题}
\begin{center}\begin{tabular}{ccccccccccccc}
\toprule
题号 & 1 & 2 & 3 & 4 & 5 & 6 & 7 & 8 & 9 & 10 & 11 & 12 \\
\midrule
答案 & D & A & C & B & A & B & C & C & D & A & B & C \\
\bottomrule
\end{tabular}\end{center}

\gkpart{填空题}
\begin{itemize}[leftmargin=2.6em]
  \item[\textbf{13.}] $20$。判断框先判 $i\leqslant n$，故 $(i,S,T)$ 依次取 $(1,1,1),(2,3,4),(3,6,10),(4,10,20)$，再判 $4>3$ 不成立，输出 $T=20$。
  \item[\textbf{14.}] $18$。可行域三个顶点 $(0,2)$、$(1,0)$、$(2,3)$ 处 $z$ 分别为 $8$、$3$、$18$。
  \item[\textbf{15.}] $12$。设 $MN$ 中点为 $Q$，则 $|AN|=2|F_1Q|$，$|BN|=2|F_2Q|$，故和为 $4a=12$。
  \item[\textbf{16.}] $-1$。$|2a+b|\leqslant 2\sqrt c$，取等时 $b=2a$、$c=4a^2$，原式为 $\dfrac2a+\dfrac1{a^2}$，当 $a=-1$ 时取最小值 $-1$。
\end{itemize}

\gkpart{解答题}
\begin{itemize}[leftmargin=2.6em]
  \item[\textbf{17.}] (1) $a=3$，$c=2$。由 $ac\cos B=2$ 得 $ac=6$，再由余弦定理得 $a^2+c^2=13$，结合 $a>c$ 解出；(2) $\dfrac{23}{27}$。$\cos C=\dfrac79$，$\sin C=\dfrac{4\sqrt2}9$，$\sin B=\dfrac{2\sqrt2}3$，代入两角差余弦公式
  \item[\textbf{18.}] (1) 有 $95\%$ 的把握。$\chi^2=\dfrac{100(60\times10-20\times10)^2}{80\times20\times70\times30}\approx4.762>3.841$；(2) $\dfrac7{10}$。$\dfrac{\mathrm C_2^0\mathrm C_3^3+\mathrm C_2^1\mathrm C_3^2}{\mathrm C_5^3}=\dfrac{1+6}{10}$
  \item[\textbf{19.}] (1) 见解析版；(2) $\dfrac12$。$DG=\dfrac{\sqrt6}2$ 为高，$S_{\triangle BCG}=\dfrac{\sqrt6}2$，故 $V=\dfrac13\cdot\dfrac{\sqrt6}2\cdot\dfrac{\sqrt6}2=\dfrac12$
  \item[\textbf{20.}] (1) $P(\sqrt2,\sqrt2)$。设切点 $(x_0,y_0)$，三角形面积 $\dfrac8{x_0y_0}$，由 $x_0^2+y_0^2=4\geqslant2x_0y_0$ 知 $x_0y_0\leqslant2$，取等时 $x_0=y_0=\sqrt2$；(2) $\dfrac{x^2}6+\dfrac{y^2}3=1$。由 $P$ 在椭圆上得 $\dfrac1{a^2}+\dfrac1{b^2}=\dfrac12$，再由弦长与面积条件解得 $b^2=3$，$a^2=6$
  \item[\textbf{21.}] (1) 见解析版；(2) 见解析版
\end{itemize}

\gkpart{选做题}
\begin{itemize}[leftmargin=2.6em]
  \item[\textbf{22.}] (1) 见解析版；(2) 见解析版
  \item[\textbf{23.}] (1) $\begin{cases}x=\cos\theta,\\y=2\sin\theta\end{cases}$（$\theta$ 为参数）。由 $x^2+\dfrac{y^2}4=1$ 直接写出；(2) $\rho=\dfrac3{4\sin\theta-2\cos\theta}$。中点 $M\left(\dfrac12,1\right)$，所求直线 $2x-4y+3=0$，代入 $x=\rho\cos\theta$，$y=\rho\sin\theta$
  \item[\textbf{24.}] (1) $M=\left[0,\dfrac43\right]$。分 $x\geqslant1$ 与 $x<1$ 去绝对值求解；(2) 见解析版
\end{itemize}
